\documentclass[a4paper,11pt]{article}
\usepackage[utf8]{inputenc}
\usepackage[T1]{fontenc}
\usepackage[french]{babel}
\usepackage{amsmath,amssymb}
\usepackage{geometry}
\usepackage{tikz}
\usetikzlibrary{arrows.meta}
\usepackage{enumitem}
\usepackage{fancyhdr}
\usepackage{tcolorbox}
\usepackage{pgfplots}
\pgfplotsset{compat=1.18}

\geometry{margin=2cm}
\setlength{\headheight}{14pt}

\pagestyle{fancy}
\fancyhf{}
\lhead{BTS CIEL -- Mathématiques}
\rhead{Évaluation -- Séries de Fourier}
\cfoot{\thepage}

\tcbset{
  exobox/.style={colback=orange!5!white, colframe=orange!60!black, fonttitle=\bfseries, title=#1},
  reponsebox/.style={colback=gray!5!white, colframe=gray!40!black, sharp corners, boxrule=0.5pt, left=4pt, right=4pt, top=4pt, bottom=4pt},
  qcmbox/.style={colback=blue!4!white, colframe=blue!50!black, fonttitle=\bfseries, title=#1}
}

% Commande case à cocher QCM
\newcommand{\choix}[1]{\hspace{0.3cm}\tikz[baseline=-0.5ex]{\draw[gray!60!black] (0,0) rectangle (0.38,0.38);}\; #1\hspace{0.6cm}}
\newcommand{\ligneqcm}[4]{\choix{#1} \choix{#2} \choix{#3} \choix{#4}}

\newcommand{\dt}{\mathrm{d}t}

\begin{document}

\begin{center}
  {\LARGE \textbf{Évaluation -- Séries de Fourier}} \\[0.2cm]
  {\Large Analyse harmonique de signaux périodiques} \\[0.2cm]
  \rule{0.6\textwidth}{0.4pt} \\[0.2cm]
  {\normalsize Durée : 1 heure \quad -- \quad Calculatrice autorisée} \\[0.1cm]
  {\normalsize Le barème est donné à titre indicatif.}
\end{center}

\vspace{0.2cm}

\noindent\textbf{Nom :} \dotfill \quad \textbf{Prénom :} \dotfill \quad \textbf{Classe :} \dotfill

\vspace{0.3cm}

\noindent\fbox{\parbox{\dimexpr\textwidth-2\fboxsep-2\fboxrule}{%
\textbf{Rappels :} Soit $f$ une fonction périodique de période $T$, de fréquence $f_0 = \dfrac{1}{T}$ et de pulsation $\omega = \dfrac{2\pi}{T}$. Son développement en série de Fourier est :
\[
  S(t) = a_0 + \sum_{n=1}^{+\infty}\!\bigl(a_n\cos(n\omega t)+b_n\sin(n\omega t)\bigr)
\]
avec \quad
$\displaystyle a_0 = \frac{1}{T}\int_0^{T} f(t)\,\dt$,\quad
$\displaystyle a_n = \frac{2}{T}\int_0^{T} f(t)\cos(n\omega t)\,\dt$,\quad
$\displaystyle b_n = \frac{2}{T}\int_0^{T} f(t)\sin(n\omega t)\,\dt$.

\smallskip
\noindent\textbf{Propriétés de symétrie :}
\begin{itemize}[leftmargin=1.5em, topsep=2pt]
  \item Si $f$ est \textbf{paire} : $b_n = 0$ pour tout $n \geq 1$.
  \item Si $f$ est \textbf{impaire} : $a_0 = 0$ et $a_n = 0$ pour tout $n \geq 1$.
\end{itemize}
}}

\vspace{0.5cm}

%% ================================================================
%% EXERCICE 1 — QCM  (5 points)
%% ================================================================
\begin{tcolorbox}[qcmbox={Exercice 1 \, -- \, QCM \hfill \normalfont\textit{(5 points)}}]
  Pour chaque question, \textbf{une seule} réponse est correcte. Cochez la case correspondante.
  \textit{Aucune justification n'est demandée.}
\end{tcolorbox}

\vspace{0.2cm}

\begin{enumerate}[label=\textbf{Q\arabic*.}, leftmargin=*, itemsep=0.5cm]

  \item La pulsation fondamentale d'un signal périodique de période $T = 5\,\mathrm{ms}$ vaut :
  \hfill \textit{(1 pt)}\\[0.25cm]
  \ligneqcm{$\omega = 200\pi\,\mathrm{rad/s}$}{$\omega = 400\pi\,\mathrm{rad/s}$}{$\omega = 200\,\mathrm{rad/s}$}{$\omega = 2000\pi\,\mathrm{rad/s}$}

  \item Pour une fonction \textbf{impaire} $f$ de période $T$, on a nécessairement :
  \hfill \textit{(1 pt)}\\[0.25cm]
  \ligneqcm{$a_0 \neq 0$ et $b_n = 0$}{$a_0 = 0$ et $a_n = 0$}{$b_n = 0$ et $a_0 = 0$}{$a_n \neq 0$ et $b_n \neq 0$}

  \item Le coefficient $a_0$ du développement en série de Fourier d'un signal représente :
  \hfill \textit{(1 pt)}\\[0.25cm]
  \ligneqcm{Son amplitude maximale}{Sa fréquence fondamentale}{Sa valeur moyenne}{Son harmonique de rang 1}

  \item On considère la fonction $f$ de période $T = 2$ définie par $f(t) = 3$ sur $[0\,;\,1[$ et $f(t) = -3$ sur $[1\,;\,2[$. Quelle est la valeur de $a_0$ ?
  \hfill \textit{(1 pt)}\\[0.25cm]
  \ligneqcm{$a_0 = 3$}{$a_0 = 1{,}5$}{$a_0 = 0$}{$a_0 = 6$}

  \item Le spectre en fréquences d'un signal périodique est constitué de raies situées :
  \hfill \textit{(1 pt)}\\[0.25cm]
  \ligneqcm{Aux fréquences $n\cdot f_0$}{Aux fréquences $f_0 / n$}{À toutes les fréquences}{À la fréquence $f_0$ uniquement}

\end{enumerate}

\vspace{0.5cm}

%% ================================================================
%% EXERCICE 2 — Signal créneaux asymétrique  (10 points)
%% ================================================================
\begin{tcolorbox}[exobox={Exercice 2 \hfill \normalfont\textit{(10 points)}}]
On considère le signal électrique $u$ périodique de période $T = 4\,\mathrm{ms}$, défini sur
$[0\,;\,T[$ par :
\[
  u(t) = \begin{cases} 6 & \text{si } 0 \leq t < 1\,\mathrm{ms} \\ 0 & \text{si } 1\,\mathrm{ms} \leq t < 4\,\mathrm{ms} \end{cases}
\]
On pose $t$ en millisecondes dans toute la suite de l'exercice.
\end{tcolorbox}

\begin{center}
  \begin{tikzpicture}
    \begin{axis}[
      axis lines=middle,
      xlabel={$t$ (ms)},
      ylabel={$u(t)$ (V)},
      xmin=-4.5, xmax=12.5,
      ymin=-1, ymax=8,
      xtick={-4,-3,-2,-1,0,1,2,3,4,5,6,7,8,9,10,11,12},
      ytick={0,2,4,6},
      width=14cm, height=4.8cm,
      grid=major,
      grid style={dashed, gray!30},
    ]
    % période [-4 ; 0[
    \addplot[blue, ultra thick, domain=-4:-3] {6};
    \addplot[blue, ultra thick, domain=-3:0]  {0};
    % période [0 ; 4[
    \addplot[blue, ultra thick, domain=0:1]   {6};
    \addplot[blue, ultra thick, domain=1:4]   {0};
    % période [4 ; 8[
    \addplot[blue, ultra thick, domain=4:5]   {6};
    \addplot[blue, ultra thick, domain=5:8]   {0};
    % période [8 ; 12[
    \addplot[blue, ultra thick, domain=8:9]   {6};
    \addplot[blue, ultra thick, domain=9:12]  {0};
    % discontinuités
    \addplot[blue, dashed, thin] coordinates {(-3,6)(-3,0)};
    \addplot[blue, dashed, thin] coordinates {(1,6)(1,0)};
    \addplot[blue, dashed, thin] coordinates {(5,6)(5,0)};
    \addplot[blue, dashed, thin] coordinates {(9,6)(9,0)};
    \end{axis}
  \end{tikzpicture}
\end{center}

\begin{enumerate}[label=\textbf{\arabic*.}]
  \item Donner la fréquence $f_0$ (en Hz) et la pulsation $\omega$ (en rad/s) de ce signal.
        \hfill \textit{(2 pts)}
  \begin{tcolorbox}[reponsebox, height=2.2cm]\end{tcolorbox}

  \item Calculer $a_0$. Interpréter physiquement ce résultat.
        \hfill \textit{(2 pts)}
  \begin{tcolorbox}[reponsebox, height=3cm]\end{tcolorbox}

  \item Calculer $a_n$ pour tout entier $n \geq 1$.
        \hfill \textit{(2 pts)}
  \begin{tcolorbox}[reponsebox, height=3.5cm]\end{tcolorbox}

  \item Calculer $b_n$ pour tout entier $n \geq 1$.
        En déduire les valeurs numériques de $b_1$, $b_2$ et $b_3$
        (valeurs arrondies au centième). \hfill \textit{(3 pts)}
  \begin{tcolorbox}[reponsebox, height=5.5cm]\end{tcolorbox}

  \item Écrire le développement en série de Fourier $S(t)$ de $u$ en ne conservant que les
        harmoniques jusqu'à $n = 3$. \hfill \textit{(1 pt)}
  \begin{tcolorbox}[reponsebox, height=2cm]\end{tcolorbox}
\end{enumerate}

\newpage

%% ================================================================
%% EXERCICE 3 — Parseval et spectre  (10 points)
%% ================================================================
\begin{tcolorbox}[exobox={Exercice 3 \hfill \normalfont\textit{(10 points)}}]
On considère un signal électrique $s$ périodique de période $T = 1\,\mathrm{ms}$,
dont le développement en série de Fourier est \textbf{donné} :
\[
  s(t) = 3 + 8\cos(\omega t) + 4\cos(2\omega t) + \frac{8}{3}\cos(3\omega t) + 2\cos(4\omega t)
\]
où $\omega = \dfrac{2\pi}{T}$.

\smallskip
On rappelle le \textbf{théorème de Parseval} pour un signal de la forme
$S(t) = a_0 + \displaystyle\sum_{n=1}^{N}a_n\cos(n\omega t)$ :
\[
  P = \frac{1}{T}\int_0^T s^2(t)\,\dt = a_0^2 + \frac{1}{2}\sum_{n=1}^{N} a_n^2
\]
$P$ est la \textbf{puissance moyenne} du signal (en $\mathrm{W}$ si $s$ est en volts sur une résistance de $1\,\Omega$).
\end{tcolorbox}

\begin{enumerate}[label=\textbf{\arabic*.}]

  \item Identifier les coefficients $a_0$, $a_1$, $a_2$, $a_3$ et $a_4$ du développement.
        \hfill \textit{(1 pt)}
  \begin{tcolorbox}[reponsebox, height=1.8cm]\end{tcolorbox}

  \item Donner la fréquence fondamentale $f_0$ (en Hz) et les fréquences des harmoniques
        de rang 1 à 4.
        \hfill \textit{(1 pt)}
  \begin{tcolorbox}[reponsebox, height=2cm]\end{tcolorbox}

  \item \textbf{Spectre unilatéral en amplitude.}
        Compléter le spectre ci-dessous en traçant les raies verticales
        à la bonne hauteur pour chaque composante fréquentielle (fondamentale + harmoniques).
        Porter les valeurs numériques au-dessus de chaque raie.
        \hfill \textit{(3 pts)}

  \begin{center}
    \begin{tikzpicture}
      \begin{axis}[
        axis lines=left,
        xlabel={Fréquence (Hz)},
        ylabel={Amplitude},
        xmin=0, xmax=5500,
        ymin=0, ymax=10,
        xtick={0,1000,2000,3000,4000,5000},
        xticklabels={$0$,$f_0$,$2f_0$,$3f_0$,$4f_0$,$5f_0$},
        ytick={0,1,2,3,4,5,6,7,8,9,10},
        width=13cm, height=6cm,
        grid=major,
        grid style={dashed, gray!20},
        ymajorgrids=true,
        tick align=outside,
      ]
      % Raie DC (a0 = 3) — tracée pour guider
      \addplot[red!70!black, ultra thick, -{Stealth[scale=1.2]}]
        coordinates {(0,0)(0,3)};
      \node[above, font=\small\bfseries, red!70!black] at (axis cs:0,3) {$3$};
      % Raies à compléter (pointillés guides)
      \foreach \x in {1000,2000,3000,4000}{
        \addplot[gray!40, dashed, thin] coordinates {(\x,0)(\x,9.5)};
      }
      \end{axis}
    \end{tikzpicture}
  \end{center}

  \item Calculer la puissance de chaque composante :
        $P_0 = a_0^2$, \; $P_n = \dfrac{a_n^2}{2}$ pour $n = 1, 2, 3, 4$.
        \hfill \textit{(2 pts)}
  \begin{tcolorbox}[reponsebox, height=3.5cm]\end{tcolorbox}

  \item Appliquer le théorème de Parseval pour calculer la \textbf{puissance totale} $P$ du signal.
        \hfill \textit{(1 pt)}
  \begin{tcolorbox}[reponsebox, height=2cm]\end{tcolorbox}

  \item Calculer le \textbf{taux de distorsion harmonique} (TDH) du signal, défini par :
        \[
          \mathrm{TDH} = \frac{\sqrt{P_2^2 + P_3^2 + P_4^2}}{P_1} \times 100\,\%
        \]
        Commenter la valeur obtenue (signal peu ou très distordu ?).
        \hfill \textit{(2 pts)}
  \begin{tcolorbox}[reponsebox, height=3.5cm]\end{tcolorbox}

\end{enumerate}

\end{document}
